\documentclass[10pt,a4paper]{article}
\usepackage[utf8]{inputenc}
\usepackage[vietnamese]{babel}
\usepackage{amsmath,amssymb,amsfonts}
\usepackage{geometry}
\geometry{a4paper, left=10mm, right=10mm, top=10mm, bottom=10mm}
\usepackage{tikz}
\usetikzlibrary{shapes,arrows.meta,positioning,calc}
\usepackage[most]{tcolorbox}
\usepackage{array,booktabs,multicol}
\usepackage{fontawesome}
\usepackage{helvet}
\renewcommand{\familydefault}{\sfdefault}
\usepackage{parskip}

\definecolor{primary}{RGB}{112, 48, 160}
\definecolor{teal}{RGB}{0, 128, 128}
\definecolor{orange}{RGB}{230, 108, 44}
\definecolor{green}{RGB}{46, 139, 87}
\definecolor{rose}{RGB}{192, 0, 0}

\pagestyle{empty}

\begin{document}

% --- HEADER BANNER ---
\begin{tcolorbox}[
  enhanced,
  colback=violet!8!white,
  colframe=violet!80!black,
  arc=2.5mm, boxrule=1pt,
  left=3mm, right=3mm, top=2mm, bottom=2mm
]
  \begin{tabular*}{\textwidth}{@{\extracolsep{\fill}} lr}
    {\small\bfseries\color{violet!85!black}\faBook\ SỔ TAY KIẾN THỨC TOÁN 2} & 
    {\small\bfseries\color{teal!85!black}\faGraduationCap\ CHƯƠNG TRÌNH GDPT 2018}
  \end{tabular*}
  \vspace{-1mm}
  \begin{center}
    \rule{0.98\textwidth}{0.6pt}
  \end{center}
  \vspace{-2mm}
  \centering
  {\LARGE\bfseries\color{violet!90!black}\faBookmark\ BÀI 63: LUYỆN TẬP CHUNG PHÉP CỘNG, PHÉP TRỪ TRONG PHẠM VI 1000}\\[1mm]
  {\itshape\footnotesize\color{black!75} Rèn luyện kĩ năng tính cộng trừ có nhớ/không nhớ và giải toán có lời văn}
\end{tcolorbox}

\vspace{-2mm}

\noindent
\begin{minipage}[t]{0.49\textwidth}\vspace{0pt}
\begin{tcolorbox}[
  enhanced,
  colback=violet!4!white,
  colframe=violet!75!black,
  arc=2mm, boxrule=0.8pt,
  left=2mm, right=2mm, top=1.8mm, bottom=1.8mm,
  title={\footnotesize\bfseries\color{white}\faKey\ 1. Khắc sâu kĩ thuật đặt tính và tính}
]
{\footnotesize
\begin{itemize}\setlength{\itemsep}{1pt}\setlength{\leftskip}{-2.5mm}
  \item \textbf{Đặt tính chuẩn xác:} Các chữ số cùng hàng phải thẳng cột với nhau (đơn vị thẳng đơn vị, chục thẳng chục, trăm thẳng trăm).
  \item \textbf{Thực hiện từ phải sang trái:}\\
  Đơn vị $\longrightarrow$ Chục $\longrightarrow$ Trăm.
  \item \textbf{Quy tắc có nhớ 1 lần:}
  \begin{itemize}\setlength{\itemsep}{0.5pt}\setlength{\leftskip}{-2mm}
    \item Phép cộng: Nhớ $1$ sang hàng bên trái liền kề (phải cộng thêm $1$).
    \item Phép trừ: Mượn $1$ ở hàng bên trái thì phải trả $1$ vào số trừ của hàng đó.
  \end{itemize}
\end{itemize}
}
\end{tcolorbox}
\end{minipage}
\hfill
\begin{minipage}[t]{0.49\textwidth}\vspace{0pt}
\begin{tcolorbox}[
  enhanced,
  colback=teal!4!white,
  colframe=teal!75!black,
  arc=2mm, boxrule=0.8pt,
  left=2mm, right=2mm, top=1.8mm, bottom=1.8mm,
  title={\footnotesize\bfseries\color{white}\faSliders\ 2. Kĩ năng tính nhẩm và giải toán}
]
{\footnotesize
\begin{itemize}\setlength{\itemsep}{1pt}\setlength{\leftskip}{-2.5mm}
  \item \textbf{Tính nhẩm số tròn trăm, tròn chục:}
  \begin{itemize}\setlength{\itemsep}{0.5pt}\setlength{\leftskip}{-2mm}
    \item $400 + 500 = 900$; \quad $800 - 300 = 500$.
    \item $360 + 20 = 380$; \quad $750 - 50 = 700$.
  \end{itemize}
  \item \textbf{Tính giá trị biểu thức không có ngoặc:} Thực hiện từ trái sang phải:
  $$a + b - c = (a + b) - c$$
  \item \textbf{Giải toán có lời văn:}\\
  Bài toán nhiều hơn ($\to$ cộng), ít hơn ($\to$ trừ).
\end{itemize}
}
\end{tcolorbox}
\end{minipage}

\vspace{-1.5mm}
\begin{tcolorbox}[
  enhanced,
  colback=orange!4!white,
  colframe=orange!80!black,
  arc=2mm, boxrule=0.8pt,
  left=2mm, right=2mm, top=1.5mm, bottom=1.5mm,
  title={\footnotesize\bfseries\color{white}\faSitemap\ 3. Sơ đồ đối chiếu quy trình cộng và trừ trong phạm vi 1000}
]
\centering
\begin{tikzpicture}[x=1.65cm, y=0.75cm]
  \tikzset{arr/.style={-{Stealth[scale=0.8]}}, every node/.append style={align=center}}
  \node[rectangle, rounded corners=2pt, draw=violet!80!black, fill=violet!15, font=\scriptsize\bfseries] (c) at (0,0.6) {PHÉP CỘNG\\$\ge 10 \implies$ Viết số đ.vị, nhớ $1$};
  \node[rectangle, rounded corners=2pt, draw=teal!80!black, fill=teal!15, font=\scriptsize\bfseries] (t) at (0,-0.6) {PHÉP TRỪ\\Không trừ được $\implies$ Mượn $1$, trả $1$};

  \node[rectangle, rounded corners=2pt, draw=green!70!black, fill=green!15, font=\footnotesize\bfseries] (r) at (3.5,0) {
    CHUNG NGUYÊN TẮC:\\
    1. Đặt thẳng cột theo hàng\\
    2. Tính từ PHẢI sang TRÁI
  };

  \draw[arr, thick] (c) -| (r);
  \draw[arr, thick] (t) -| (r);
\end{tikzpicture}
\end{tcolorbox}

\vspace{-1.5mm}
\noindent
\begin{minipage}[t]{0.49\textwidth}\vspace{0pt}
\begin{tcolorbox}[
  enhanced,
  colback=green!4!white,
  colframe=green!60!black,
  arc=2mm, boxrule=0.8pt,
  left=2mm, right=2mm, top=1.8mm, bottom=1.8mm,
  title={\footnotesize\bfseries\color{white}\faCheckCircle\ 4. Bài tập thực hành mẫu}
]
{\footnotesize
\textbf{1. Đặt tính rồi tính:}\\
- $428 + 315 = \mathbf{743}$ ($8+5=13$, nhớ $1$).\\
- $652 - 238 = \mathbf{414}$ ($12-8=4$, nhớ $1$).\\
\textbf{2. Tính giá trị biểu thức:}\\
$300 + 450 - 150 = 750 - 150 = \mathbf{600}$.\\
\textbf{3. Bài toán lời văn:}\\
Đội Một trồng được $345$ cây, đội Hai trồng nhiều hơn đội Một $50$ cây. Hỏi đội Hai trồng được bao nhiêu cây?\\
\textit{Giải:} Đội Hai trồng được: $345 + 50 = \mathbf{395}$ (cây).
}
\end{tcolorbox}
\end{minipage}
\hfill
\begin{minipage}[t]{0.49\textwidth}\vspace{0pt}
\begin{tcolorbox}[
  enhanced,
  colback=rose!4!white,
  colframe=rose!70!black,
  arc=2mm, boxrule=0.8pt,
  left=2mm, right=2mm, top=1.8mm, bottom=1.8mm,
  title={\footnotesize\bfseries\color{white}\faExclamationTriangle\ 5. Lời dặn dò quan trọng}
]
{\footnotesize
\begin{itemize}\setlength{\itemsep}{1.5pt}\setlength{\leftskip}{-2.5mm}
  \item \textbf{Thói quen kiểm tra:} Sau khi tính xong, lấy kết quả cộng/trừ ngược lại để thử lại ngay.
  \item Đọc kĩ đề toán lời văn để không nhầm giữa "nhiều hơn" ($+$) và "ít hơn" ($-$).
  \item Khi số hạng chỉ có 2 chữ số (ví dụ $45$), nhớ để trống vị trí hàng trăm ở phía trước.
\end{itemize}
}
\end{tcolorbox}
\end{minipage}


\end{document}
